\documentclass[11pt,letterpaper]{report}

\usepackage[margin=1in]{geometry}
\usepackage[T1]{fontenc}
\usepackage{lmodern}
\usepackage{microtype}
\usepackage{amsmath,amssymb}
\usepackage{booktabs}
\usepackage{tabularx}
\usepackage{array}
\usepackage{enumitem}
\usepackage{xcolor}
\usepackage{listings}

\usepackage[hidelinks]{hyperref}
\usepackage{xurl}
\hypersetup{
  pdftitle={HPINNS User Manual},
  pdfauthor={HPINNS},
  pdfsubject={How to train classical, hybrid-quantum, and full-quantum shallow-water PINNs}
}

\setlist[itemize]{leftmargin=1.4em,itemsep=0.25em,topsep=0.4em}
\setlist[enumerate]{leftmargin=1.5em,itemsep=0.25em,topsep=0.4em}
\setlist[description]{style=nextline,leftmargin=1.2em,itemsep=0.45em,topsep=0.4em}

\lstdefinestyle{py}{
  language=Python,
  basicstyle=\ttfamily\footnotesize,
  breaklines=true,
  breakatwhitespace=false,
  frame=single,
  framesep=5pt,
  xleftmargin=0.4em,
  framexleftmargin=0.4em,
  columns=fullflexible,
  keepspaces=true,
  showstringspaces=false,
  tabsize=4,
  aboveskip=0.85em,
  belowskip=0.7em,
  keywordstyle=\color{blue!60!black}\bfseries,
  commentstyle=\color{black!55},
  stringstyle=\color{teal!50!black}
}
\lstdefinestyle{sh}{
  language=bash,
  basicstyle=\ttfamily\footnotesize,
  breaklines=true,
  frame=single,
  framesep=5pt,
  xleftmargin=0.4em,
  framexleftmargin=0.4em,
  columns=fullflexible,
  keepspaces=true,
  showstringspaces=false,
  aboveskip=0.85em,
  belowskip=0.7em,
  keywordstyle=\color{blue!60!black}\bfseries,
  commentstyle=\color{black!55}
}
\lstset{style=py}

\DeclareUrlCommand{\file}{\urlstyle{tt}}
\DeclareUrlCommand{\code}{\urlstyle{tt}}
\setlength{\emergencystretch}{2em}

\title{%
  HPINNS User Manual\\[0.4em]
  \large Classical, hybrid-quantum, and full-quantum\\
  shallow-water physics-informed neural networks}
\author{HPINNS}
\date{November 23, 2026}

\begin{document}
\maketitle

\begin{abstract}
\noindent
HPINNS trains a physics-informed neural network for the shallow-water equations.
One trainer serves three model families.
The classical family gives each field its own multilayer perceptron.
The hybrid-quantum family keeps that perceptron and adds a residual quantum circuit on every tower.
The full-quantum family has no hidden perceptron: each field is a linear readout of a variational circuit applied to the scaled coordinates.
This manual explains how to install the package, choose a family, set a benchmark, run a fixed model, search several families with Optuna, and read the files a run writes.
Worked input settings are given for each family and for a joint three-family search.
\end{abstract}

\tableofcontents

\chapter{Introduction}
\label{ch:intro}

HPINNS is a Python package for shallow-water physics-informed neural networks.
The network represents water depth and velocity.
Training minimizes the mean square of the shallow-water residuals at a fixed grid of collocation points, together with a boundary penalty when constraints are soft.
When the selected benchmark has an exact solution, a later analysis step compares the trained field with that solution.

The package name on disk is \file{swe_pinn}.
The program you run is \file{run_pinn.py}.
Every setting that changes from one experiment to the next lives in \file{user_input.py}.

\section{What one run does}

A call to \code{python run_pinn.py} does the following.

\begin{enumerate}
\item Read \code{USER_PINN_INPUTS} from \file{user_input.py}.
\item Build the benchmark named by \code{problem_name}.
  The physical constants for that benchmark come from the matching dataclass at the bottom of the same file.
\item If a search is requested, draw trials until the completed-trial cap is reached.
  Each trial builds one network, trains it with Adam, and scores the same loss on a finer lattice.
\item Rebuild the best completed trial, or the fixed hyperparameters if the search was skipped, and train that model with Adam for \code{final_adam_epochs}.
\item Optionally continue from those weights with L-BFGS-B.
\item Write weights, a hyperparameter summary, and NumPy archives of the loss history and the predicted fields.
\end{enumerate}

Two flags split that sequence so several processes can share one study.

\begin{itemize}
\item \code{python run_pinn.py --search-only --worker-id N} runs trials only.
  The sampler seed is \(42+N\).
  The process does not train the final model.
\item \code{python run_pinn.py --finalize} loads the best completed trial already stored in the study and runs the final Adam stage, the optional L-BFGS-B stage, and the archives.
  It adds no trials.
\end{itemize}

The two flags cannot be combined.
\file{run_cores.sh} starts several search workers, waits for them, and then runs \code{--finalize} once.

\section{The three families}

Every benchmark can be trained with any of the three families.
The family is a property of the network, not of the equation.

\begin{description}
\item[\code{classical}]
Each field has its own tanh multilayer perceptron.
There is no circuit.
Width and depth set the size of each tower.

\item[\code{hybrid_quantum}]
Each field keeps that perceptron.
After the last hidden layer the tower adds a residual quantum correction,
\[
z \leftarrow z + \operatorname{decode}\!\big(\operatorname{circuit}(\operatorname{encode}(z))\big).
\]
The correction is added at full strength on every forward pass after an optional classical warmup.
A hybrid specification with zero qubits is stored and trained as classical.

\item[\code{full_quantum}]
There is no hidden perceptron and no residual skip.
Each field is
\[
\operatorname{readout}\!\big(\operatorname{circuit}(\operatorname{encode}(\boldsymbol{\xi}))\big),
\]
where \(\boldsymbol{\xi}\) is the pair of coordinates scaled onto the unit interval.
Width and depth are ignored.
The qubit count and the layer count must both be at least one.
\end{description}

Chapter~\ref{ch:families} specifies the circuit, the two embedding layouts, and the parameter counts.
Chapter~\ref{ch:examples} gives a complete fixed-architecture run for each family.
Chapter~\ref{ch:search} shows how to put all three families in one Optuna study.

\section{What this manual assumes}

You edit \file{user_input.py}, run the solver from the repository root, and read the text files and archives under \file{analysis/}.
You do not need to edit the trainer to change the family, the benchmark, the constraint mode, or the search.
A new benchmark is added by registering a problem builder, which is described briefly in Section~\ref{sec:new-problem}.

\chapter{Installation and layout}
\label{ch:install}

\section{Requirements}

The trainer uses 64-bit floating point for every trainable tensor and every PDE derivative.
\file{swe_pinn/physics/quantum.py} enables JAX double precision when the package is imported, so JAX is imported even for a classical run.

Install the third-party packages from the repository root:

\begin{lstlisting}[style=sh]
pip install tensorflow optuna scipy matplotlib numpy pennylane jax jaxlib pandas
\end{lstlisting}

\begin{center}
\begin{tabular}{@{}lp{9.2cm}@{}}
\toprule
Package & Role \\
\midrule
TensorFlow / Keras &
Dense layers, Adam, \code{GradientTape} derivatives, saved weights \\
NumPy &
Collocation coordinates, exact-solution arrays, archives \\
Optuna &
Study, TPE sampler, SQLite storage, completed-trial cap \\
SciPy &
Optional L-BFGS-B; root finding for the dam-break and bump references \\
PennyLane &
Circuit definition and the \code{default.qubit} statevector device \\
JAX &
Batched just-in-time circuit and its vector--Jacobian products \\
Matplotlib &
Figures written by \file{analysis/analyze_study.py} \\
pandas &
Trial tables in the analysis script. The trainer does not import it \\
\bottomrule
\end{tabular}
\end{center}

PennyLane is imported only when a circuit is constructed.
A classical-only run still imports JAX, because the quantum module is imported with the package, and it does not construct a circuit.
A hybrid or full-quantum run needs both PennyLane and JAX.
The circuit is compiled with the JAX interface and differentiated with \code{diff_method="backprop"}.

Python~3.10 or newer matches the type hints in \file{user_input.py}.
The commands below use \code{python}; on some systems the interpreter is \code{python3}.
\file{run_cores.sh} calls \code{python3}.

\section{Check that the package imports}

From the repository root:

\begin{lstlisting}[style=sh]
python -c "from swe_pinn.physics.common import MODEL_FAMILIES; print(MODEL_FAMILIES)"
\end{lstlisting}

The expected line is

\begin{lstlisting}[style=sh]
('classical', 'hybrid_quantum', 'full_quantum')
\end{lstlisting}

A missing PennyLane installation does not fail this import.
It fails later, on the first hybrid or full-quantum model, with a message that names \code{pip install pennylane}.

\section{Repository layout}

Run every command from the repository root, the directory that contains \file{run_pinn.py} and \file{user_input.py}.

\begin{center}
\begin{tabularx}{\textwidth}{@{}lX@{}}
\toprule
Path & Role \\
\midrule
\file{user_input.py} &
The active input file. The solver reads \code{USER_PINN_INPUTS} and the physical dataclass selected by \code{problem_name}. \\
\file{user_input.py.*} &
Ready-made copies of that file for particular benchmarks and family sets. Copy one over \file{user_input.py} to use it. \\
\file{run_pinn.py} &
Process entry point. \\
\file{run_cores.sh} &
\(N\) search workers, then one final training. \\
\file{swe_pinn/cli.py} &
Argument parser for \code{--search-only}, \code{--worker-id}, and \code{--finalize}. \\
\file{swe_pinn/core/solver.py} &
Loss, Adam, L-BFGS-B, Optuna, and archives. \\
\file{swe_pinn/physics/common.py} &
The \code{PINN} network, collocation grids, and the problem dataclass. \\
\file{swe_pinn/physics/quantum.py} &
PennyLane circuit and the TensorFlow bridge. \\
\file{swe_pinn/problems/} &
Benchmark builders and the registry in \file{__init__.py}. \\
\file{swe_pinn/plots/jcp.py} &
Plot helpers. The solver itself writes NumPy archives; the study figures are made by the analysis script. \\
\file{analysis/analyze_study.py} &
Tradeoff report, Optuna figures, and the analytic comparison. \\
\file{analysis/<database>/} &
One folder per SQLite file: the database, \file{data/}, and \file{figures/}. \\
\file{logs/} &
Worker logs and the shared trial-table counter used by a parallel search. \\
\bottomrule
\end{tabularx}
\end{center}

A relative SQLite URL such as \code{sqlite:///optuna_tide.db} is stored at \file{analysis/optuna_tide/optuna_tide.db}.
The folder name is the database file name without \file{.db}.
Weights, \file{best_hyperparameters.txt}, and the NumPy archives go in \file{analysis/optuna_tide/data/}, whether or not Optuna actually ran.
Point \code{storage_url} at a new file before a smoke test so a short run does not replace the weights of a finished study.
Absolute SQLite URLs are left unchanged.

\section{First command}

After the packages import, the shortest confirmation that the entry point is wired up is the help text:

\begin{lstlisting}[style=sh]
python run_pinn.py --help
\end{lstlisting}

That command does not train a model.
The first training run is Example~\ref{ex:classical} in Chapter~\ref{ch:examples}.
Give it a new \code{storage_url} and a modest \code{final_adam_epochs} before you launch a multi-hour study.

\chapter{The three model families}
\label{ch:families}

The network class is \code{PINN} in \file{swe_pinn/physics/common.py}.
Every family accepts two scaled coordinates and returns one scalar per field, concatenated along the last axis.
One-dimensional benchmarks return \((h,u)\) or \((\eta,u)\).
The three steady square benchmarks return \((h,u,v)\).
Each field has its own tower.
There is no shared trunk.

Coordinates are scaled onto \([0,1]\) before they enter the network.
Dense-layer kernels use Glorot normal initialization and zero biases.
Hidden layers use \(\tanh\).
Output layers are linear.
Circuit angles are drawn from a normal distribution with mean zero and standard deviation \(0.1\).
\code{build_model} resets the NumPy and TensorFlow seeds to 42 before constructing the network, so a given architecture starts from the same weights in every trial and in the final training.

The legal names are \code{classical}, \code{hybrid_quantum}, and \code{full_quantum}.
Any other string raises \code{ValueError} and lists those three names.

\section{Classical}
\label{sec:classical}

Set
\begin{lstlisting}
model_family = "classical"
\end{lstlisting}
Each tower is a stack of \code{depth} hidden layers of width \code{width}, followed by a linear unit.
The circuit is not built.
\code{n_qubits}, \code{q_layers}, and \code{embedding} are ignored, and the stored embedding is \code{once}.

For one tower with two inputs and depth \(d\ge 1\), the trainable count is
\begin{equation}
\label{eq:classical}
P_{\mathrm{C}}
= 3w + (d-1)(w^{2}+w) + (w+1),
\end{equation}
where \(w\) is the width.
The first term is the map from the two coordinates into the first hidden layer.
The last term is the linear readout.
An \(m\)-field model has \(m P_{\mathrm{C}}\) parameters, all of them classical.
The banner prints \code{quantum_params=0}.

The shipped search ranges use \(d\ge 2\).
Equation~\eqref{eq:classical} is the count for that regime.
A depth of~1 is a single hidden layer: \(P_{\mathrm{C}}=4w+1\).

\section{Hybrid quantum}
\label{sec:hybrid}

Set
\begin{lstlisting}
model_family = "hybrid_quantum"
n_qubits_value = 4
q_layers_value = 2
embedding_value = "reupload"
\end{lstlisting}
The classical towers are unchanged.
After the last hidden layer of each tower the forward pass adds
\begin{equation}
\label{eq:hybrid}
z \leftarrow z +
\operatorname{decode}\!\big(
\operatorname{circuit}(\operatorname{encode}(z))
\big).
\end{equation}
The corrected state then enters the linear output layer.
Each tower has its own encoder, circuit, and decoder.
The correction is not multiplied by a trainable gate.

The encoder maps the \(w\)-dimensional hidden state to one feature per qubit.
The decoder is a \(\tanh\) dense map from the Pauli-\(Z\) expectations back to width \(w\).
With the \code{once} embedding the encoder uses \(\tanh\).
With \code{reupload} the encoder is linear, so a coordinate weight can grow past one radian.

A hybrid value of \code{n_qubits} equal to 0 is rewritten as \code{classical} before training, and the circuit is not built.
Keep \code{n_qubits_min} at 1 or more when the name is being searched, so a quantum trial actually contains a circuit.
\code{q_layers} must be at least 1 whenever the circuit is built.

For one or more qubits the count per tower is
\begin{equation}
\label{eq:hybrid-count}
P_{\mathrm{H}}
= P_{\mathrm{C}}
+ (wn_{q}+n_{q})
+ (n_{q}w+w)
+ 3n_{q}n_{L}.
\end{equation}
The three added terms are the encoder, the decoder, and the circuit rotations.
The banner splits them: encoder and decoder are classical, and the \(3n_{q}n_{L}\) angles are quantum.
An \(m\)-field model multiplies \(P_{\mathrm{H}}\) by \(m\).

\subsection{Classical warmup}

\code{hybrid_classical_warmup_fraction} is the fraction of Adam epochs, in a trial and again in the final stage, during which a hybrid tower skips the circuit.
The quantum variables still exist.
Their gradients are replaced by zeros so Adam's variable list does not change when the circuit turns on.
A fraction of 0 keeps the circuit on from epoch 1.
The fraction is clamped to the epoch budget.
The shipped \file{user_input.py} sets the fraction to 0.

Warmup applies only to \code{hybrid_quantum}.
A full-quantum model has no classical trunk to train first, and the fraction does not change that family.
After a hybrid trial the circuit is switched back on before the score lattice is evaluated.

\section{Full quantum}
\label{sec:full}

Set
\begin{lstlisting}
model_family = "full_quantum"
n_qubits_value = 4
q_layers_value = 2
embedding_value = "reupload"
\end{lstlisting}
Width and depth are stored as 0 and are not sampled.
Both \code{n_qubits} and \code{q_layers} must be at least 1.
A smaller value raises \code{ValueError} and names the field to raise.

Let \(\boldsymbol{\xi}\in\mathbb{R}^{2}\) be the two coordinates after scaling onto the unit interval.
For the tide, the dam break, and the bump, those coordinates are scaled space and time.
For the steady square problems they are scaled \(x\) and \(y\).
Each tower is
\begin{equation}
\label{eq:full}
\operatorname{readout}\!\big(
\operatorname{circuit}(\operatorname{encode}(\boldsymbol{\xi}))
\big).
\end{equation}
The encoder maps the two coordinates to \(n_{q}\) features.
The readout maps the \(n_{q}\) expectations to one scalar.
Both maps are classical dense layers.
There is no residual skip around the circuit.

The count per tower is
\begin{equation}
\label{eq:full-count}
P_{\mathrm{F}}
= 4n_{q} + 1 + 3n_{q}n_{L}.
\end{equation}
The term \(4n_{q}+1\) is the encoder (\(3n_{q}\) weights) plus the linear readout (\(n_{q}+1\)).
The banner reports those weights as classical and the \(3n_{q}n_{L}\) angles as quantum.
An \(m\)-field model has \(m P_{\mathrm{F}}\) parameters.
The family name means that the tower has no hidden perceptron.
The encoder and the readout are still classical trainable parameters.

\section{The circuit}
\label{sec:circuit}

Both quantum families use \code{QuantumLayer}.
The trainable tensor \(\theta\) has shape \code{(q_layers, n_qubits, 3)}.
That is \(3 n_{q} n_{L}\) angles per tower.

Two layouts are legal: \code{once} and \code{reupload}.
A classical model does not use either one.

\begin{description}
\item[\code{once}]
The encoder uses \(\tanh\).
\code{AngleEmbedding} places those features on the qubits with \(Y\) rotations, once, before the variational layers.
Each layer is then an input-independent \code{Rot} on every qubit, followed by the CNOT ring.

\item[\code{reupload}]
The encoder is linear.
The same features are embedded again, as \(Y\) rotations, at the start of every variational layer, and the trainable \code{Rot} follows that embedding.
\end{description}

After the rotations, a layer with more than one qubit applies CNOT gates from qubit \(i\) to \(i+1\) and a closing CNOT from the last qubit to the first.
A one-qubit layer is only the rotation.
The circuit returns the Pauli-\(Z\) expectation of every qubit.

The device is PennyLane \code{default.qubit}.
JAX vectorizes the circuit over the collocation batch with one shared set of angles, then just-in-time compiles the circuit and its derivatives.
TensorFlow still owns \(\theta\) and the tape that forms the PDE residual.
A custom gradient returns the circuit's vector--Jacobian product with respect to the features and the angles.
The feature product is differentiated again, which is the second derivative a residual needs when a slope of the network updates the angles.
The first quantum model in a process pays the compilation cost.
Later epochs of the same architecture reuse the compiled circuit.

Fix \code{embedding_value} to give every quantum tower the same layout.
Put \code{"embedding"} in \code{tune_variables} only when a study should compare layouts, and list the candidates in \code{embedding_options}.
The shipped studies fix \code{reupload} and leave \code{embedding} out of the search.

\section{Parameter counts you can check}
\label{sec:counts}

The banner printed at the start of Adam training reports classical, quantum, and total trainable parameters.
Table~\ref{tab:counts} is the count from Equations~\eqref{eq:classical}--\eqref{eq:full-count}.
A run that disagrees with this table is using a different width, depth, qubit count, layer count, or number of fields.

\begin{table}[ht]
\centering
\caption{Trainable-parameter counts for the formulas in this chapter.
Two-field models are the tide, the dam break, and the bump.
Three-field models are the steady square problems.}
\label{tab:counts}
\begin{tabular}{@{}llrrrrrr@{}}
\toprule
Family & Fields & \(w\) & \(d\) & \(n_{q}\) & \(n_{L}\) & Quantum & Total \\
\midrule
classical & 2 & 22 & 6 & 0 & 0 & 0 & 5238 \\
classical & 3 & 5 & 7 & 0 & 0 & 0 & 603 \\
classical & 3 & 1 & 2 & 0 & 0 & 0 & 21 \\
hybrid & 2 & 22 & 6 & 4 & 2 & 48 & 5690 \\
hybrid & 3 & 1 & 2 & 1 & 1 & 9 & 42 \\
full quantum & 2 & --- & --- & 4 & 2 & 48 & 82 \\
full quantum & 3 & --- & --- & 1 & 1 & 9 & 24 \\
\bottomrule
\end{tabular}
\end{table}

The three-field hybrid row is the lower edge of the parameter window in the shipped eddy study: width~1, depth~2, one qubit, and one layer.
The three-field classical row with width~5 and depth~7 is the upper edge of that same window.
Section~\ref{sec:window} explains why the window is set from those two edges.

\section{What each family is for}

Use the classical family when you want a baseline perceptron, when PennyLane is not installed, or when a search should vary only width and depth.

Use the hybrid family when the perceptron should carry the smooth part of the field and the circuit should supply a residual correction.
Warmup is available only for this family.
Width and depth still matter, and they are sampled together with the circuit size when those names are listed.

Use the full-quantum family when the field itself should be a readout of the circuit.
Width and depth do not change the model.
The capacity is the qubit count, the layer count, and the embedding.
The encoder and readout remain classical, so the banner's classical count is not zero.

All three families share the loss, the collocation grids, the Adam loop, the optional L-BFGS-B stage, and the archive format.
A comparison across families is a comparison of those networks on that shared trainer.
Chapter~\ref{ch:search} is the procedure for putting them in one study.

\chapter{Training, constraints, and the loss}
\label{ch:training}

\section{The scalar loss}

One loss is used for Adam, for L-BFGS-B, and for the Optuna score.
Let \(r_{m}\) be the mass residual and let \(r_{p,j}\) be the momentum residuals returned by the benchmark.
One-dimensional problems return one momentum residual.
The steady square problems return two.
The logged names are mass and momentum for every benchmark, including the linear tide.
The PDE pieces are
\begin{equation}
L_{m} = \operatorname{mean}(r_{m}^{2}),
\qquad
L_{p} = \sum_{j} \operatorname{mean}(r_{p,j}^{2}).
\end{equation}
Derivatives are TensorFlow reverse-mode derivatives through the output map, including the hard-constraint transformation.
There is no test function and no integration by parts.

Logarithmic parameters \(\lambda_{m}\) and \(\lambda_{p}\) (\code{log_w_mass} and \code{log_w_mom}) set the relative equation weights:
\begin{equation}
\label{eq:weights}
(w_{m}, w_{p})
=
\frac{\big(e^{\lambda_{m}}, e^{\lambda_{p}}\big)}
{\tfrac12\big(e^{\lambda_{m}}+e^{\lambda_{p}}\big)}.
\end{equation}
Equal logarithms give \((w_{m},w_{p})=(1,1)\), and the two weights always sum to 2.
These weights multiply only the PDE terms.
The soft penalty has its own coefficients on the physical dataclass.
The scalar loss is
\begin{equation}
\label{eq:loss}
L = w_{m} L_{m} + w_{p} L_{p} + L_{\mathrm{soft}}.
\end{equation}
In hard mode, \(L_{\mathrm{soft}}=0\).
The IC and BC columns in the epoch table are the pieces inside \(L_{\mathrm{soft}}\).
They are already included in the total.
They are not added a second time.

\section{Hard and soft constraints}
\label{sec:constraints}

\code{constraints} is \code{hard} or \code{soft}.
Any other string raises \code{ValueError}.

Soft mode returns the network output unchanged and adds a penalty that pulls selected traces onto prescribed data.
The penalty weights live on the physical dataclass (\code{soft_w_*} fields), not in the PDE log-weights.

Hard mode builds the output so the prescribed traces are satisfied for any finite network output.
The penalty is omitted.
The interior value is that boundary construction plus a factor that vanishes on the constrained boundary, multiplied by the network.
The construction is different for each benchmark, and it is specified in Chapter~\ref{ch:benchmarks}.
In particular, hard mode on the solid-body rotation inserts the exact solution in the interior.
Hard mode on the vortex and on the eddy lattice lifts the boundary traces only.

Every built-in benchmark defines a soft penalty, so both modes are legal on all six problems.
A custom problem that omits \code{soft_constraint_loss_fn} cannot be run with \code{constraints="soft"}.

\section{Collocation}

Interior points are a fixed tensor product, including both endpoints of each axis.
Points are not redrawn between epochs.
The flatten order is every station of the second axis at the first \(x\), then every station of the second axis at the next \(x\).
For the steady square problems the second axis is \(y\).
The input fields are still named \code{trial_nt}, \code{score_nt}, \code{final_adam_nt}, and so on.

\begin{center}
\footnotesize
\begin{tabular}{@{}ll@{}}
\toprule
Setting & Lattice \\
\midrule
\code{trial_nx}, \code{trial_nt} & Adam inside each Optuna trial \\
\code{score_nx}, \code{score_nt} & Objective lattice, after training \\
\code{final_adam_nx}, \code{final_adam_nt} & Final Adam stage \\
\code{final_lbfgs_nx}, \code{final_lbfgs_nt} & L-BFGS-B, when enabled \\
\code{nx}, \code{nt} & Evaluation mesh for archives and the analytic comparison \\
\bottomrule
\end{tabular}
\end{center}

Both station counts must be at least 2.
The score lattice does not take another Adam step.
It is the same strong-form loss, including the soft penalty, on a grid that is usually finer than the training grid.
The shipped eddy study trains on \(48\times 48\) and scores on \(96\times 96\).
The figure mesh, often \(201\times 201\), is not a training grid.
Optuna trials do not write a model-versus-analytic figure.
That comparison is written once, by \file{analysis/analyze_study.py}, after a model has been trained.

\section{Adam}

Adam uses a constant learning rate \code{lr}.
The optimizer rescales each weight from the gradient moments.
The scalar itself does not decay.
Exponential decay is not exposed in \file{user_input.py}; a decaying rate left the soft dam-break loss on a plateau in the reference setup, and the trainer keeps the constant rate.

The epoch table has six columns: epoch, total loss, mass loss, momentum loss, IC data, and BC data.
\code{show_epoch_output} and \code{epoch_output_every} control the final stage.
\code{show_trial_epoch_output} and \code{trial_epoch_output_every} control the table inside each trial.
A row is printed at epoch~1 and every \(N\)th epoch after that.
The header is printed once per stage.

\code{trial_epochs} is the budget used to rank a trial.
\code{final_adam_epochs} is a separate, usually longer, budget for the model that is kept.
The final stage is not itself an Optuna trial.

\section{L-BFGS-B}

\code{use_lbfgs} continues from the Adam weights with SciPy L-BFGS-B on the same trainable variables and the same total loss.
\code{final_lbfgs_max_iters} is an iteration cap.
\code{final_lbfgs_maxfun} is a separate cap on loss and gradient evaluations.
\code{final_lbfgs_ftol} stops on a small relative decrease of the loss.
\code{final_lbfgs_gtol} stops on a small projected gradient.
The shipped \file{user_input.py} sets \code{use_lbfgs = False}, which keeps the Adam weights and still writes the summary and the archives.
The summary then records that L-BFGS-B was skipped.

\section{How a trial is scored}
\label{sec:objective}

After \code{trial_epochs} of Adam, the trial is ranked on \(L\) evaluated on the score lattice.
The value stored as \code{raw_objective} is that loss, including the soft penalty.
Optuna may minimize a transformed number.
\code{parameter_objective_mode} selects the transformation, and \code{optimize_total_parameters} is the master switch.
Parameter counts are stored on the trial in every mode.

\begin{description}
\item[\code{raw_loss}]
The objective is \(L\).
This is also the value returned when \code{optimize_total_parameters} is false, for every mode.

\item[\code{loss_multiplier}]
The objective is
\[
L \left(1 + \alpha \frac{N}{N_{\mathrm{ref}}}\right),
\]
where \(N\) is the trainable-parameter count, \(\alpha\) is \code{total_params_objective_weight}, and \(N_{\mathrm{ref}}\) is \code{total_params_reference}.
A weight of 0 returns \(L\).

\item[\code{target_error_min_params}]
The objective is
\[
\frac{N}{N_{\mathrm{ref}}}
+ W \max\!\left(0,\, \frac{L - L_{\star}}{L_{\star}}\right),
\]
where \(L_{\star}\) is \code{target_raw_objective} and \(W\) is \code{target_error_constraint_weight}.
The excess term is zero once \(L \le L_{\star}\).
This mode requires \(L_{\star}>0\).
Set \(L_{\star}\) from a finished classical baseline on the same problem before using the mode.
The multiplier weight \(\alpha\) is logged and does not enter this formula.
\end{description}

The score is a residual loss.
It is not an \(L^{2}\) error against the exact depth or velocity.
Field errors are computed on the evaluation mesh after the final model is trained.

\section{Reading the trial table}

Each finished trial prints one row to the terminal, including when the worker's standard output has been redirected to a log file.
The row contains the trial number, family, width, depth, qubits, layers, an embedding column when the layout can vary, the learning rate and log-weights when those values can differ from the fixed defaults, the classical, quantum, and total parameter counts, the mass, momentum, IC, and BC pieces of the score, the Optuna objective, the wall-clock minutes, and the count of completed trials against the cap.

A draw whose parameter count falls outside \code{total_params_min} or \code{total_params_max} is discarded before training.
The terminal says \code{Trial N discarded}, the trial is pruned, and it does not count toward \code{n_trials}.
A trial that raises is printed as \code{FAILED} together with a traceback in the process log.
Failed trials do not count as completed.

\chapter{Built-in benchmarks}
\label{ch:benchmarks}

\code{problem_name} selects the builder.
The physical numbers come from the dataclass instantiated at the bottom of \file{user_input.py}.
The solver does not mix parameters across benchmarks.
Changing \code{g}, a domain limit, or a boundary weight means editing that dataclass, not \code{USER_PINN_INPUTS}.

All six benchmarks define an exact or reference solution, a soft penalty, and a hard output map.
Table~\ref{tab:problems} is the index.
The second coordinate of the steady problems is \(y\).

\begin{table}[ht]
\centering
\caption{Built-in benchmarks. The field count is the number of network towers.}
\label{tab:problems}
\begin{tabular}{@{}llll@{}}
\toprule
\code{problem_name} & Fields & Second axis & Default idea \\
\midrule
\code{tidal_test} & 2 & time & linear tide, exact solution \\
\code{dam_break} & 2 & time & wet-bed Stoker problem \\
\code{transcritical_bump} & 2 & dummy time & steady jump over a bump \\
\code{steady_rotation} & 3 & \(y\) & solid-body rotation \\
\code{steady_vortex} & 3 & \(y\) & gradient-wind vortex \\
\code{steady_eddies} & 3 & \(y\) & Fourier eddies on a carved bed \\
\bottomrule
\end{tabular}
\end{table}

\section{Linear tidal test}

\code{problem_name = "tidal_test"} builds a linear tide in a flat basin of length \(L\), undisturbed depth \(H\), and linear friction \(r\).
The free-surface amplitude at the forced end is \(A\), and the period is \(T_{\mathrm{tide}}\) (the default is 12.42~hours).
The trainer works in nondimensional \(x\in[0,1]\) and \(\tau\) from 0 to \code{tau_max}, with \(\tau\) in units of the tidal period.
The default window is three tidal periods.
The comparison archives convert elevation back by \(A\), velocity by \(A L/(H T_{\mathrm{tide}})\), \(x\) to kilometers, and time to hours.

The residuals are the linear continuity and momentum balances,
\[
\eta_{\tau} + u_{x} = 0,
\qquad
u_{\tau} + \alpha \eta_{x} + \beta u = 0,
\]
with \(\alpha = g H T_{\mathrm{tide}}^{2}/L^{2}\) and \(\beta = r T_{\mathrm{tide}}\).
The exact solution is the complex wave number that satisfies those equations with the basin boundary data.

Hard mode matches that exact solution on \(\tau=0\), on \(x=0\), and on \(x=1\).
The interior is the transfinite blend of those three traces plus \(\tau\, x\,(1-x)\) times the network.
The collocation loss is the PDE residual.
The exact interior solution is not pasted into the domain.

Soft mode penalizes the network against the exact initial elevation and velocity, and against the exact elevation and velocity on both ends.
The default weights, used when the dataclass does not override them, are 20 on each initial field and 10 on each boundary field.
Sample counts default to 2048 initial points and 1024 points on each end, drawn uniformly at random on each penalty evaluation.

\section{Wet-bed dam break}

\code{problem_name = "dam_break"} is the scaled Stoker problem used by the reference dam-break study:
\(g=9.81\,\mathrm{m\,s^{-2}}\), left depth \(0.5\,\mathrm{m}\), right depth \(0.25\,\mathrm{m}\), still water, \(x\in[0,1]\,\mathrm{m}\), \(t\in[0,0.2]\,\mathrm{s}\), dam at \(x=0.5\,\mathrm{m}\).
The reference solution is the wet-bed Riemann solution.

The residuals are the one-dimensional conservation laws
\[
h_{t} + (hu)_{x} = 0,
\qquad
(hu)_{t} + \big(hu^{2} + \tfrac12 g h^{2}\big)_{x} = 0.
\]

\code{transition_width = 0} keeps the sharp initial step.
A positive width replaces the step by a \(\tanh\) jump of that scale, in meters.

Hard mode sets
\[
h = h_{0}(x) + \hat{t}\,\hat{x}\,(1-\hat{x})\, h_{\mathrm{net}},
\qquad
u = u_{0} + \hat{t}\,\hat{x}\,(1-\hat{x})\, u_{\mathrm{net}},
\]
where hats denote coordinates scaled to \([0,1]\) and \(h_{0}\) is the initial step.
At \(t=0\) and on both walls the network contribution vanishes, so the depth on each wall stays at the initial depth of that wall and the velocity stays at \(u_{0}\).

Soft mode penalizes the initial step and the two walls separately.
The default loss is
\[
L_{m} + L_{p} + 50\,\mathrm{MSE}(h_{\mathrm{IC}}) + 50\,\mathrm{MSE}(u_{\mathrm{IC}})
+ 10\,\mathrm{MSE}(h_{\mathrm{BC}}) + 10\,\mathrm{MSE}(u_{\mathrm{BC}}).
\]
The default 41 initial samples are packed toward the dam.
Each wall has 10 samples, uniform in time.
The network sees time in the first column and \(x\) in the second.
The other benchmarks pass \(x\) first.
You do not set that order; it is part of the dam-break builder.

\section{Transcritical bump}

\code{problem_name = "transcritical_bump"} is steady, frictionless flow over a parabolic bump with a stationary hydraulic jump.
The channel is \(x\in[0,25]\,\mathrm{m}\).
The bed is \(z = 0.2 - 0.05(x-10)^{2}\) for \(8<x<12\) and zero elsewhere.
The upstream unit discharge is \(q=0.18\,\mathrm{m^{2}\,s^{-1}}\) and the downstream depth is \(h=0.33\,\mathrm{m}\).
The reference solution follows from constant discharge and the specific-energy relation on each smooth branch.
The jump is placed where the Rankine--Hugoniot momentum condition changes sign.

The residuals are the steady balances with bed slope and no friction,
\[
(hu)_{x} = 0,
\qquad
\big(hu^{2} + \tfrac12 g h^{2}\big)_{x} + g h z_{x} = 0.
\]
The trainer still carries a time coordinate, on the dummy interval \([0,1]\).
The network is always evaluated at \(t=0\), so the learned field does not depend on time.

Hard mode enforces the mixed physical data exactly:
downstream depth is \(h(L)=h_{\mathrm{down}}\) for any network output, and upstream discharge is \(h(0)u(0)=q_{\mathrm{in}}\).
Upstream depth and downstream velocity remain free.
Soft mode penalizes those two conditions, each with default weight 50, and does not prescribe both \(h\) and \(u\) at each end.

\section{Steady solid-body rotation}

\code{problem_name = "steady_rotation"} is a steady square flow,
\[
u = -\omega y, \qquad v = \omega x,
\qquad
h = h_{\mathrm{center}} + \frac{\omega^{2}(x^{2}+y^{2})}{2g},
\]
on a flat frictionless bed.
The default square is \([-1,1]\times[-1,1]\), with \(\omega=1\,\mathrm{s^{-1}}\) and \(h_{\mathrm{center}}=1\,\mathrm{m}\).
The residuals are the steady conservative shallow-water equations with no friction and no bed slope.
There is no time axis.
\code{trial_nx} and \code{trial_nt} are the \(x\) and \(y\) station counts.

Hard mode returns the exact field plus a bubble \(s_{x}(1-s_{x})s_{y}(1-s_{y})\) times the network, where \(s_{x}\) and \(s_{y}\) are the coordinates scaled to \([0,1]\).
The bubble is zero on the boundary, so the boundary values are exact.
A zero network is already the exact solution everywhere, because the exact field is part of the output.
This benchmark is a useful wiring check.
It is a weak test of whether a family can represent an unknown interior.

Soft mode penalizes \(h\), \(u\), and \(v\) on the four edges against the exact traces.
The default edge has 32 stations, and each field weight defaults to 50.

\section{Steady vortex}

\code{problem_name = "steady_vortex"} is an isolated vortex in gradient-wind balance on a flat bed.
The azimuthal speed is \(\omega r \exp(-r^{2}/(2\sigma^{2}))\), the radial velocity is zero, and the free surface satisfies
\[
\frac{v_{\theta}^{2}}{r} + f v_{\theta} = g \frac{\partial h}{\partial r}.
\]
The default Coriolis parameter is \(f=2\,\mathrm{s^{-1}}\), a rotating-table value rather than the planetary one.
Set \(f=0\) on \code{SteadyVortexParameters} for a non-rotating bucket vortex.
The default core rate is \(\omega=4\,\mathrm{s^{-1}}\), the speed peaks at \(r=\sigma=0.45\,\mathrm{m}\), and the central depth is \(0.40\,\mathrm{m}\).
The square is \([-1.5,1.5]\times[-1.5,1.5]\), several vortex radii across, so the boundary velocity is a small fraction of the peak.
Continuity is satisfied identically.
There is no manufactured body force.

Hard mode lifts the boundary traces with a Coons patch and adds the same boundary-vanishing bubble times the network.
The exact interior solution is not inserted.
Soft mode penalizes all three fields on the four edges against the exact boundary values, with the same default station count and weights as the rotation problem.

\section{Steady eddy lattice}

\code{problem_name = "steady_eddies"} is the benchmark selected by the shipped \file{user_input.py}.
The stream function is a degenerate Fourier pair
\[
\psi = A \big[
\sin(m\pi x/L)\sin(n\pi y/L)
+ \sin(n\pi x/L)\sin(m\pi y/L)
\big],
\]
with velocity \((\psi_{y},-\psi_{x})\) and depth \(H_{0}+\gamma\psi\).
Both mode pairs share the eigenvalue \(\mu=(m^{2}+n^{2})\pi^{2}/L^{2}\).
The free surface that balances advection is
\[
\eta = C - \frac{|\mathbf{u}|^{2}}{2g} - \frac{\mu\psi^{2}}{2g},
\]
and the bed is \(z=\eta-h\).
Continuity and momentum hold identically.
There is no manufactured body force and no shock.
The walls are slip walls: depth equals \(H_{0}\) and the normal velocity vanishes.
The defaults are \(A=0.015\,\mathrm{m^{2}\,s^{-1}}\), \(H_{0}=0.80\,\mathrm{m}\), \(\gamma=4\), modes \((m,n)=(7,4)\), Bernoulli constant \(C=1\), and the unit square.
Keep \(C\) large enough that the bed stays below the free surface.
Mode~7 needs more than twenty stations across the square; the shipped trial lattice uses 48.

Hard mode is the Coons lift of the boundary traces plus the bubble times the network.
The exact interior is not inserted.
Soft mode enforces the slip condition only: depth equals \(h_{\mathrm{mean}}\) on the walls, and the normal velocity vanishes.
Tangential speed is left free so that penalty does not dominate the Optuna score.
\code{soft_w_u} and \code{soft_w_v} are averaged into the normal-velocity weight.

\section{Choosing a constraint mode}

Use hard mode when the boundary traces should be exact for every set of weights, and the loss should be the PDE residual alone.
That is the shipped choice for the tide and for the eddy lattice.

Use soft mode when the boundary data are mixed, or when you want the penalty weights to be part of the experiment.
The dam-break reference and the bump preset use soft mode.
The bump's hard map already enforces the mixed discharge and depth conditions, so hard mode is also a valid bump run.
Give the bump a new \code{study_name} if you change the constraint mode of an existing study; the trials are not comparable across modes.

Use a new database whenever you change \code{problem_name}.
Trials from the tide and trials from the dam break must not share a study.
The loss values are residuals of different equations.

\section{Adding a benchmark}
\label{sec:new-problem}

A new case is a builder that returns a \code{Problem} from \file{swe_pinn/physics/common.py}: a name, a \code{Domain1D}, a residual function, an output function, and, when you want soft mode or an analytic comparison, a soft-penalty function and an exact solution.
Register the name in \code{build_problem} in \file{swe_pinn/problems/__init__.py}, and add a dataclass to \file{user_input.py}.
The residual returns two tensors for a one-dimensional problem or three tensors for a steady \((x,y)\) problem.
Set \code{n_outputs} to 2 or 3 to match.
The trainer, the three families, and the Optuna loop stay as they are.
Point \code{storage_url} and \code{study_name} at a new study before the first trial.

\chapter{The input file}
\label{ch:config}

The solver imports \file{user_input.py} and reads the object \code{USER_PINN_INPUTS}.
A name listed in \code{tune_variables} is sampled by Optuna on each trial, inside the matching range.
Every other setting stays at its \code{*_value} field.
\code{model_family} is the exception among the categorical settings: when it is tuned, Optuna draws from \code{model_family_options}.
\code{embedding} is drawn from \code{embedding_options}, and only after the trial is still \code{hybrid_quantum} or \code{full_quantum}.

If \code{use_optuna} is false, if \code{tune_variables} is empty, or if \code{n_trials} is 0, no study is created and the fixed values are trained directly.
Artifacts still go to the folder named by \code{storage_url}.

\section{Run identity}

\begin{description}
\item[\code{problem_name}]
One of the six names in Table~\ref{tab:problems}.

\item[\code{constraints}]
\code{hard} or \code{soft}. See Section~\ref{sec:constraints}.

\item[\code{model_family}]
The family used when \code{"model_family"} is not in \code{tune_variables}.

\item[\code{model_family_options}]
The families a tuned study may draw.
Remove a name to keep that family out of the search.
An existing study freezes this set.
Change \code{study_name} before adding or removing a family, or Optuna rejects the new categorical distribution.

\item[\code{storage_url}]
SQLite URL of the study.
A relative file is stored in \file{analysis/<database name>/}.
A new URL starts a new database.
Trials are appended when the file already exists.

\item[\code{study_name}]
Name of the study inside that database.
Reloading the same name continues the same search, including its frozen distributions.
Use a new name when the family list, the embedding list, or an integer search range changes.

\item[\code{use_optuna}]
Master switch for the search.
\end{description}

\section{Fixed architecture}

These fields are used when the corresponding name is absent from \code{tune_variables}.
They are also the values a search falls back to for any name it is not allowed to sample.

\begin{description}
\item[\code{width_value}, \code{depth_value}]
Hidden width and hidden depth of each classical tower.
Ignored for \code{full_quantum}.
The dam-break reference search preferred a width near 22.

\item[\code{lr_value}]
Adam learning rate for every family.
The shipped studies leave \code{"lr"} out of \code{tune_variables}, so this value is the rate.
Adding \code{"lr"} to the list samples it for every family, log-uniform between \code{lr_min} and \code{lr_max}.

\item[\code{log_w_mass_value}, \code{log_w_mom_value}]
The logarithms in Equation~\eqref{eq:weights}.
Zeros give PDE weights \((1,1)\).

\item[\code{n_qubits_value}, \code{q_layers_value}]
Circuit size for a fixed hybrid or full-quantum model.
A hybrid zero is stored as classical.
A full-quantum value below 1 aborts the run.

\item[\code{embedding_value}]
\code{once} or \code{reupload}.
Classical trials ignore it.
The shipped file sets \code{reupload}.

\item[\code{embedding_options}]
Layouts a search may draw when \code{"embedding"} is tuned.
The shipped tuple contains only \code{reupload}.
This set is frozen with the study in the same way as \code{model_family_options}.
\end{description}

\section{What a search is allowed to sample}

\code{tune_variables} accepts these names and no others:

\begin{center}
\begin{tabular}{@{}ll@{}}
\toprule
Name & Distribution when it is listed \\
\midrule
\code{model_family} & categorical over \code{model_family_options} \\
\code{width} & integers from the width box, stride \code{width_step} \\
\code{depth} & integers from the depth box, stride 1 \\
\code{lr} & log-uniform between \code{lr_min} and \code{lr_max} \\
\code{log_w_mass}, \code{log_w_mom} & uniform on the matching min/max \\
\code{n_qubits}, \code{q_layers} & integers on the matching min/max \\
\code{embedding} & categorical over \code{embedding_options} \\
\bottomrule
\end{tabular}
\end{center}

Sampling is conditional on the family drawn for that trial.

\begin{center}
\begin{tabular}{@{}p{3.3cm}p{5.4cm}p{5.4cm}@{}}
\toprule
Family & Sampled when listed & Left at the fixed value \\
\midrule
\code{classical} &
width, depth, learning rate, log-weights &
qubits, layers, embedding \\
\code{hybrid_quantum} &
those, plus qubits, layers, and embedding &
--- \\
\code{full_quantum} &
learning rate, log-weights, qubits, layers, embedding &
width and depth, which are stored as 0 \\
\bottomrule
\end{tabular}
\end{center}

A hybrid draw with zero qubits is rewritten as classical and then follows the classical rules.
With \code{n_qubits_min} at 1 or more, that rewrite does not occur.

\section{Family-specific ranges}

When \code{use_family_specific_search_ranges} is true, a classical trial uses \code{classical_width_*} and \code{classical_depth_*}, and a hybrid trial uses the \code{hybrid_*} bounds.
Full quantum does not sample width or depth in either case.
When the flag is false, both families share \code{width_min}, \code{width_max}, \code{width_step}, \code{depth_min}, and \code{depth_max}.

The shipped eddy study sets the two boxes to the same trunk, width 1 through 5 and depth 2 through 7.
The comparison is then the residual circuit, rather than a different range of perceptrons.
A hybrid whose circuit pushes the total over \code{total_params_max} is discarded.

Inclusive bounds and the width stride are the Optuna suggestions
\[
w = w_{\min},\; w_{\min}+w_{\mathrm{step}},\; \ldots,\; w_{\max}.
\]
Depth is every integer from the minimum through the maximum.
The learning-rate floor in the shipped file is \(10^{-4}\).
A draw near \(2\times 10^{-5}\) stayed on the soft-loss plateau for a whole trial in the reference dam-break setup.
The shipped log-weight box is \([-1,1]\) for each equation.

\section{Parameter window}
\label{sec:window}

\code{total_params_min} and \code{total_params_max} discard a draw before training.
Either bound may be \code{None}, which turns that side off.
A positive bound is inclusive in the sense that a count strictly below the minimum, or strictly above the maximum, is pruned.

The shipped eddy window is 42 to 603 parameters.
Forty-two is the smallest three-tower hybrid in the search box (width 1, depth 2, one qubit, one layer).
Six hundred three is the largest three-tower classical model in that box (width 5, depth 7).
The overlap is the range of sizes that all three families can occupy, so a family is not rewarded only for being smaller or only for being larger.
Recompute the window when you change the number of fields or the width and depth boxes.
Table~\ref{tab:counts} is the arithmetic.
A two-field problem with the same width and depth boxes has a different total, and the eddy numbers 42 and 603 do not transfer to it.

Discarded draws do not count toward \code{n_trials}.
When either bound is set, the search keeps drawing until \code{n_trials} trials have completed inside the window.
When both bounds are unset, a single process stops after \code{n_trials} attempts, and a failed attempt consumes one of them.

\section{Search budget and sampler}

\begin{description}
\item[\code{n_trials}]
Completed-trial cap.
One \code{python run_pinn.py} stops at this many new completed trials, counting trials already stored in the study.
\file{run_cores.sh} shares the same cap across workers.

\item[\code{optuna_n_startup_trials}]
Random trials before the TPE sampler starts concentrating.
The sampler is Optuna's \code{TPESampler} with \code{multivariate=True} and \code{group=True}, which is the setting used for conditional parameters such as qubits and layers.
The seed is 42 in a single process and \(42+N\) for worker \(N\).
The shipped value is 2000, above the completed cap of 500, so the whole search stays in the random phase.
Each draw then picks a family with equal probability, and the three finished counts stay close.
Lower the startup count below \code{n_trials} when you want TPE to exploit the trials it has already seen.

\item[\code{trial_epochs}, \code{trial_nx}, \code{trial_nt}]
Adam budget and training lattice inside a trial.

\item[\code{score_nx}, \code{score_nt}]
Lattice for the objective. The shipped choice is twice the trial lattice on each axis.
\end{description}

\section{Final stage}

\begin{description}
\item[\code{final_adam_epochs}, \code{final_adam_nx}, \code{final_adam_nt}]
Budget and lattice for the model that is kept.
This run starts from the same seed-42 initialization as a trial.
It does not resume the trial weights.

\item[\code{show_epoch_output}, \code{epoch_output_every}]
Final-stage table.

\item[\code{show_trial_epoch_output}, \code{trial_epoch_output_every}]
Per-trial table.
The shipped file prints every trial epoch, which is a large log.
Raise the stride, or set the flag false, when the log files dominate the disk.

\item[\code{use_lbfgs} and the \code{final_lbfgs_*} fields]
Section on L-BFGS-B in Chapter~\ref{ch:training}.

\item[\code{nx}, \code{nt}]
Evaluation mesh for the archives. Not a training grid.
\end{description}

\section{Objective fields}

\code{optimize_total_parameters}, \code{parameter_objective_mode}, \code{target_raw_objective}, \code{target_error_constraint_weight}, \code{total_params_objective_weight}, and \code{total_params_reference} are defined in Section~\ref{sec:objective}.
The shipped eddy file minimizes the raw validation loss and still records parameter counts.
\code{total_params_reference} must be positive.
It is the parameter count that counts as one in the two parameter-aware formulas.
The shipped reference is \(10^{5}\).

\section{Presets}

These files are complete copies of the input file.
Copy one over \file{user_input.py} from the repository root.
That replaces the active file, so keep a copy of any local edits first.

\begin{center}
\begin{tabular}{@{}lll@{}}
\toprule
File & Benchmark & Families in the search \\
\midrule
\file{user_input.py} & eddy lattice, hard & classical, hybrid, full \\
\file{user_input.py.tide} & tide, hard & classical, hybrid, full \\
\file{user_input.py.tide_2fam} & tide, hard & classical, hybrid \\
\file{user_input.py.tide_3fam} & tide, hard, reupload & classical, hybrid, full \\
\file{user_input.py.damBreak} & dam break, soft & classical, hybrid, full \\
\file{user_input.py.dam_break_3fam} & dam break, soft, reupload & classical, hybrid, full \\
\file{user_input.py.bump} & bump, soft & classical, hybrid, full \\
\file{user_input.py.steady_rotation_2fam} & rotation, soft & classical, hybrid \\
\file{user_input.py.steady_rotation_3fam} & rotation, soft & classical, hybrid, full \\
\bottomrule
\end{tabular}
\end{center}

\begin{lstlisting}[style=sh]
cp user_input.py user_input.py.local-backup
cp user_input.py.tide_3fam user_input.py
python run_pinn.py
\end{lstlisting}

Each preset names its own database and study.
Loading a preset and then editing \code{model_family_options} without changing \code{study_name} will collide with trials already stored under that name.
Chapter~\ref{ch:examples} uses fresh database names for that reason.
The presets themselves are the production studies, with hundreds of trials and thousands of epochs.
Read \code{n_trials} and \code{final_adam_epochs} in the file you copied before you launch it.

\chapter{Examples for each family}
\label{ch:examples}

Each example is a fixed architecture: Optuna is off, and one family is trained.
Edit the assignments inside \code{UserPINNInputs} in \file{user_input.py}.
Do not paste a second copy of the class at the bottom of the file.
Leave the physical dataclasses alone unless you intend to change the benchmark.

Every example points \code{storage_url} at a new database.
The archives, including \file{best_hyperparameters.txt} and the weights, are written under \file{analysis/<database name>/data/} even though no study is created.
Reusing a production database name replaces that study's saved weights.

The epoch counts below are short enough to confirm that the family builds, trains, and writes its files.
A production run uses the budgets in the preset you started from: several thousand trial epochs and tens of thousands of final Adam epochs.
Raise \code{final_adam_epochs}, and the collocation counts, after the short run has written a summary whose family and parameter counts match the banner you expect.

Run each example from the repository root:

\begin{lstlisting}[style=sh]
python run_pinn.py
\end{lstlisting}

The process prints the problem name, the constraint mode, the family, an empty tune set, and whether L-BFGS-B is on.
It then prints \code{Optuna search skipped. Using fixed values from user_input.py} and the final Adam banner.

\section{Classical dam break}
\label{ex:classical}

This run is a soft wet-bed dam break with two perceptrons and no circuit.
Width 22 and depth 6 are the scale of the reference dam-break trunk.
Two towers at that size have 5238 trainable parameters, all classical.

\begin{lstlisting}
problem_name = "dam_break"
constraints = "soft"
model_family = "classical"
use_optuna = False
tune_variables = ()

width_value = 22
depth_value = 6
lr_value = 1.0e-3
log_w_mass_value = 0.0
log_w_mom_value = 0.0

storage_url = "sqlite:///optuna_manual_classical_dambreak.db"
study_name = "manual_classical_dambreak"

use_lbfgs = False
final_adam_epochs = 200
final_adam_nx = 41
final_adam_nt = 41
show_epoch_output = True
epoch_output_every = 50
nx = 101
nt = 101
\end{lstlisting}

The banner should contain
\code{model_family=classical}, \code{width=22}, \code{depth=6},
\code{n_qubits=0}, \code{q_layers=0},
\code{quantum_params=0}, and \code{total_params=5238}.
The IC and BC columns of the epoch table are nonzero, because the constraint mode is soft.
They fall when the network approaches the initial step and the wall data.
The mass and momentum columns are the conservation residuals.

After the run, confirm the summary:

\begin{lstlisting}[style=sh]
sed -n '1,20p' analysis/optuna_manual_classical_dambreak/data/best_hyperparameters.txt
\end{lstlisting}

The same directory holds \file{best_optuna_model.weights.h5}, \file{best_optuna_model_config.json}, \file{dam_break_losses_data.npz}, and the elevation, velocity, error, and boundary archives.
PennyLane is not required for this example.
JAX is imported with the package and is not asked to build a circuit.

\section{Hybrid-quantum tide}
\label{ex:hybrid}

This run is a hard linear tide.
Each field keeps a perceptron of width 16 and depth 4, and adds a two-layer, four-qubit reupload circuit.
Warmup is off, so the residual circuit is present from epoch 1.
Hard mode makes the IC and BC columns identically zero; the total loss is the PDE residual.

\begin{lstlisting}
problem_name = "tidal_test"
constraints = "hard"
model_family = "hybrid_quantum"
use_optuna = False
tune_variables = ()

width_value = 16
depth_value = 4
n_qubits_value = 4
q_layers_value = 2
embedding_value = "reupload"
hybrid_classical_warmup_fraction = 0.0

lr_value = 1.0e-3
log_w_mass_value = 0.0
log_w_mom_value = 0.0

storage_url = "sqlite:///optuna_manual_hybrid_tide.db"
study_name = "manual_hybrid_tide"

use_lbfgs = False
final_adam_epochs = 100
final_adam_nx = 32
final_adam_nt = 32
show_epoch_output = True
epoch_output_every = 25
nx = 101
nt = 101
\end{lstlisting}

Per tower, Equation~\eqref{eq:classical} with \(w=16\) and \(d=4\) gives
\(P_{\mathrm{C}}=3\cdot 16 + 3\cdot(256+16) + 17 = 881\).
Equation~\eqref{eq:hybrid-count} then adds the encoder \(16\cdot 4+4=68\), the decoder \(4\cdot 16+16=80\), and the angles \(3\cdot 4\cdot 2=24\), so one tower has 1053 parameters and the model has 2106.
Of those, \(2\cdot 24=48\) are quantum angles and 2058 are classical weights.
The banner should report those three numbers, \code{embedding=reupload}, and four qubits with two layers.

The first epoch is slower than a classical step of the same grid, because JAX compiles the circuit.
Later epochs reuse that compilation.
If PennyLane is missing, the run stops in \code{QuantumLayer} with the install instruction from \file{swe_pinn/physics/quantum.py}.

To see warmup, set \code{hybrid_classical_warmup_fraction = 0.25} and \code{final_adam_epochs = 40}.
The log should say that epochs 1 through 10 skip the QNode and that the residual is enabled from epoch 11.
The quantum parameters are still listed in the banner during the warmup epochs.

\section{Full-quantum vortex}
\label{ex:full}

This run is a hard gradient-wind vortex.
There is no hidden perceptron.
Width and depth may be left at whatever the file already contains; the trainer stores them as 0.
Four qubits and three reupload layers give, by Equation~\eqref{eq:full-count},
\[
P_{\mathrm{F}} = 4\cdot 4 + 1 + 3\cdot 4\cdot 3 = 53
\]
parameters per tower.
Three towers give 159 parameters, of which \(3\cdot 3\cdot 4\cdot 3 = 108\) are circuit angles and 51 are the classical encoders and readouts.

\begin{lstlisting}
problem_name = "steady_vortex"
constraints = "hard"
model_family = "full_quantum"
use_optuna = False
tune_variables = ()

n_qubits_value = 4
q_layers_value = 3
embedding_value = "reupload"

lr_value = 1.0e-3
log_w_mass_value = 0.0
log_w_mom_value = 0.0

storage_url = "sqlite:///optuna_manual_full_vortex.db"
study_name = "manual_full_vortex"

use_lbfgs = False
final_adam_epochs = 100
final_adam_nx = 24
final_adam_nt = 24
show_epoch_output = True
epoch_output_every = 25
nx = 81
nt = 81
\end{lstlisting}

The banner should contain \code{model_family=full_quantum}, \code{width=0}, \code{depth=0}, \code{n_qubits=4}, \code{q_layers=3}, \code{embedding=reupload}, \code{quantum_params=108}, and \code{total_params=159}.
The momentum column is the sum of the two momentum mean squares.
The archives are \file{steady_vortex_elevation_compare_data.npz}, \file{steady_vortex_velocity_u_compare_data.npz}, \file{steady_vortex_velocity_v_compare_data.npz}, and \file{steady_vortex_error_data.npz}, plus the loss archive.
Hard mode on this problem lifts the boundary.
A small loss means the PDE residual is small.
It does not, by itself, mean the interior matches the exact vortex; that comparison is the analytic report in Chapter~\ref{ch:outputs}.

A full-quantum run with \code{n_qubits_value = 0} or \code{q_layers_value = 0} stops immediately and names the field that must be at least 1.
Setting \code{hybrid_classical_warmup_fraction} has no effect on this family.

\section{Same problem, three fixed models}
\label{ex:three-fixed}

A fair look at the three families on one equation uses the same benchmark, the same constraint mode, the same learning rate, the same grids, and the same epoch budget, and writes each family to its own database.
The eddy lattice is the shipped three-family problem.
The blocks below share the benchmark, the hard constraint, the learning rate, the grids, and the epoch budget.
The classical and hybrid blocks also share the trunk.
The full-quantum block sets the circuit, and width and depth do not enter that model.
Each block has its own database.
Hard constraints, a zero warmup, unit PDE weights, and the learning rate \(10^{-3}\) are shared.

Classical, 603 parameters (width 5, depth 7, three towers):

\begin{lstlisting}
problem_name = "steady_eddies"
constraints = "hard"
model_family = "classical"
use_optuna = False
tune_variables = ()
width_value = 5
depth_value = 7
lr_value = 1.0e-3
log_w_mass_value = 0.0
log_w_mom_value = 0.0
hybrid_classical_warmup_fraction = 0.0
storage_url = "sqlite:///optuna_manual_eddies_classical.db"
study_name = "manual_eddies_classical"
use_lbfgs = False
final_adam_epochs = 200
final_adam_nx = 32
final_adam_nt = 32
show_epoch_output = True
epoch_output_every = 50
\end{lstlisting}

Hybrid, the same trunk plus one qubit and one reupload layer.
One tower has \(P_{\mathrm{C}}=201\), plus an encoder of 6 weights, a decoder of 10, and 3 angles, so \(P_{\mathrm{H}}=220\).
Three towers have 660 parameters, of which 9 are quantum angles.
That model is larger than the shipped search window: the window's 42-parameter hybrid is the other corner, width~1 and depth~2 with the same one-qubit circuit.
This example keeps the trunk lined up with the classical run, so the extra 57 parameters are only the circuit path.

\begin{lstlisting}
problem_name = "steady_eddies"
constraints = "hard"
model_family = "hybrid_quantum"
use_optuna = False
tune_variables = ()
width_value = 5
depth_value = 7
n_qubits_value = 1
q_layers_value = 1
embedding_value = "reupload"
lr_value = 1.0e-3
log_w_mass_value = 0.0
log_w_mom_value = 0.0
hybrid_classical_warmup_fraction = 0.0
storage_url = "sqlite:///optuna_manual_eddies_hybrid.db"
study_name = "manual_eddies_hybrid"
use_lbfgs = False
final_adam_epochs = 200
final_adam_nx = 32
final_adam_nt = 32
show_epoch_output = True
epoch_output_every = 50
\end{lstlisting}

The hybrid banner should report 660 total parameters and 9 quantum parameters.
The trunk matches the classical model.
The extra 57 parameters are the three encoders, the three decoders, and the nine angles.
Full quantum cannot use width 5 and depth 7.
Match it on the circuit instead, here one qubit and one layer (24 parameters, three towers):

\begin{lstlisting}
problem_name = "steady_eddies"
constraints = "hard"
model_family = "full_quantum"
use_optuna = False
tune_variables = ()
n_qubits_value = 1
q_layers_value = 1
embedding_value = "reupload"
lr_value = 1.0e-3
log_w_mass_value = 0.0
log_w_mom_value = 0.0
storage_url = "sqlite:///optuna_manual_eddies_full.db"
study_name = "manual_eddies_full"
use_lbfgs = False
final_adam_epochs = 200
final_adam_nx = 32
final_adam_nt = 32
show_epoch_output = True
epoch_output_every = 50
\end{lstlisting}

Compare the final total-loss lines and, after Chapter~\ref{ch:outputs}, the relative \(L_{2}\) errors.
Do not compare these short runs with a finished preset study.
The preset trains much longer, on a \(48\times 48\) trial lattice and a \(64\times 64\) final lattice.

\section{What the config file records}

\file{best_optuna_model_config.json} is the hyperparameter record of the model that was trained.
A classical run stores \code{model_family} as \code{classical}, qubits and layers as 0, and the width and depth you set.
A hybrid run stores the width, depth, qubits, layers, and embedding.
A full-quantum run stores width 0, depth 0, and the circuit settings.
The analysis script reloads weights only when this file matches the trial it is about to compare.
A mismatch causes it to train the architecture again.

\chapter{Searching across families}
\label{ch:search}

A search draws a new architecture for each trial, trains it for \code{trial_epochs}, and records the score.
The final Adam run happens after the cap is reached, on the single best completed trial, and it starts from a fresh initialization.

\section{Three families, one study}
\label{ex:three-search}

This is the shipped eddy experiment, reduced to a handful of trials so you can see one row from each family.
Copy the ranges from the active \file{user_input.py} and change the budget, the database, and the study name.
The production file uses 500 completed trials, 4000 epochs per trial, and 45000 final Adam epochs.

\begin{lstlisting}
problem_name = "steady_eddies"
constraints = "hard"
model_family = "hybrid_quantum"  # used only when the family is not tuned

model_family_options = (
    "classical",
    "hybrid_quantum",
    "full_quantum",
)
embedding_value = "reupload"
embedding_options = ("reupload",)

use_optuna = True
tune_variables = (
    "model_family",
    "width",
    "depth",
    "n_qubits",
    "q_layers",
)

use_family_specific_search_ranges = True
classical_width_min = 1
classical_width_max = 5
classical_width_step = 1
classical_depth_min = 2
classical_depth_max = 7
hybrid_width_min = 1
hybrid_width_max = 5
hybrid_width_step = 1
hybrid_depth_min = 2
hybrid_depth_max = 7

n_qubits_min = 1
n_qubits_max = 8
q_layers_min = 1
q_layers_max = 8

# Three-tower overlap of the box above. See Table of parameter counts.
total_params_min = 42
total_params_max = 603

lr_value = 1.0e-3
log_w_mass_value = 0.0
log_w_mom_value = 0.0
hybrid_classical_warmup_fraction = 0.0
parameter_objective_mode = "raw_loss"
optimize_total_parameters = True

n_trials = 9
optuna_n_startup_trials = 2000
trial_epochs = 50
trial_nx = 24
trial_nt = 24
score_nx = 32
score_nt = 32

final_adam_epochs = 100
final_adam_nx = 32
final_adam_nt = 32
show_trial_epoch_output = False
show_epoch_output = True
epoch_output_every = 25
use_lbfgs = False

storage_url = "sqlite:///optuna_manual_eddies_three_family.db"
study_name = "manual_eddies_three_family"
\end{lstlisting}

Launch it with one process:

\begin{lstlisting}[style=sh]
python run_pinn.py
\end{lstlisting}

What you should see:

\begin{itemize}
\item A line naming the joint search, the three families, the tune set, \code{n_startup_trials=2000}, and the completed-trial cap.
\item A line saying the search runs until 9 completed trials, each with at least 42 and at most 603 parameters.
\item Discard messages for draws outside that window.
  Those trials are absent from the completed count.
\item One table row per completed trial.
  Classical rows show qubits and layers as 0.
  Full-quantum rows show width and depth as 0.
  Hybrid rows show a positive width, depth, qubit count, and layer count, and the embedding \code{reupload}.
\item Because the startup count is above the cap, families are drawn with equal probability.
  Nine trials will not split exactly three and three and three, but a single family should not take the whole table.
\item After the ninth completion, a \code{Best trial params} listing and a final Adam banner for whichever completed trial had the lowest objective.
\end{itemize}

The learning rate is absent from \code{tune_variables}, so every family uses \(10^{-3}\).
The log-weights are absent, so every family uses PDE weights \((1,1)\).
Width and depth are sampled for classical and hybrid trials and skipped for full-quantum trials.
Qubits and layers are sampled for the two quantum families and skipped for classical trials.

The database is
\file{analysis/optuna_manual_eddies_three_family/optuna_manual_eddies_three_family.db}.
Rerunning the same command continues that study.
If nine trials are already complete, the process adds none and goes straight to final training.
Raise \code{n_trials} to append more, or change both \code{storage_url} and \code{study_name} to start over.

\section{One family, circuit size only}

To tune the circuit of a hybrid model and hold the trunk fixed, omit \code{model_family}, \code{width}, and \code{depth} from the tune list.

\begin{lstlisting}
problem_name = "tidal_test"
constraints = "hard"
model_family = "hybrid_quantum"
model_family_options = ("hybrid_quantum",)
use_optuna = True
tune_variables = ("n_qubits", "q_layers")

width_value = 16
depth_value = 4
lr_value = 1.0e-3
log_w_mass_value = 0.0
log_w_mom_value = 0.0
embedding_value = "reupload"
n_qubits_min = 1
n_qubits_max = 4
q_layers_min = 1
q_layers_max = 3

total_params_min = None
total_params_max = None
n_trials = 6
optuna_n_startup_trials = 6
trial_epochs = 40
trial_nx = 24
trial_nt = 24
score_nx = 32
score_nt = 32
final_adam_epochs = 80
final_adam_nx = 32
final_adam_nt = 32
show_trial_epoch_output = False
use_lbfgs = False

storage_url = "sqlite:///optuna_manual_hybrid_circuit.db"
study_name = "manual_hybrid_circuit"
\end{lstlisting}

Every completed row should say \code{hybrid_quantum}, width 16, and depth 4.
Qubits and layers should vary inside the boxes above.
\code{model_family_options} is unused while the family is not tuned, but a one-element tuple documents the intent and keeps a later edit from accidentally opening the other families inside this study name.
If you later add \code{"model_family"} to \code{tune_variables} on this same \code{study_name}, change the study name first.

The same pattern with \code{model_family = "full_quantum"} and the same tune list searches only the circuit.
Width and depth in the file are ignored.
A classical circuit search is empty: drop \code{n_qubits} and \code{q_layers}, and list \code{"width"} and \code{"depth"} instead.

\section{Classical width and depth}

\begin{lstlisting}
problem_name = "dam_break"
constraints = "soft"
model_family = "classical"
use_optuna = True
tune_variables = ("width", "depth")
use_family_specific_search_ranges = True
classical_width_min = 8
classical_width_max = 24
classical_width_step = 8
classical_depth_min = 2
classical_depth_max = 4

n_trials = 4
optuna_n_startup_trials = 4
trial_epochs = 30
trial_nx = 21
trial_nt = 21
score_nx = 41
score_nt = 41
final_adam_epochs = 60
show_trial_epoch_output = False
use_lbfgs = False
total_params_min = None
total_params_max = None

storage_url = "sqlite:///optuna_manual_classical_search.db"
study_name = "manual_classical_search"
\end{lstlisting}

Widths fall on 8, 16, and 24.
Depths fall on 2, 3, and 4.
The table's qubit and layer columns stay at 0.

\section{Learning rate and loss weights}

Add \code{"lr"} to \code{tune_variables} to sample the Adam rate log-uniformly between \code{lr_min} and \code{lr_max}, for whichever family the trial drew.
Add \code{"log_w_mass"} or \code{"log_w_mom"} to sample the logarithms in Equation~\eqref{eq:weights}.
Sampling either logarithm prints both log-weight columns.
Leaving both at zero and out of the tune list hides those columns, because the weights are then identically \((1,1)\).

Add \code{"embedding"} only when \code{embedding_options} contains more than one layout.
A one-element tuple makes every quantum trial use that layout, and the embedding column stays hidden unless the fixed value is something other than \code{once} or the name is being tuned.
The shipped reupload studies show the column because \code{embedding_value} is \code{reupload}.

\section{Parameter-aware objectives}

To minimize the raw loss, keep the shipped mode:

\begin{lstlisting}
parameter_objective_mode = "raw_loss"
optimize_total_parameters = True
\end{lstlisting}

\code{optimize_total_parameters = False} also returns the raw loss, for any mode.
Counts are still stored.

To charge a model for its size, set

\begin{lstlisting}
optimize_total_parameters = True
parameter_objective_mode = "loss_multiplier"
total_params_objective_weight = 1.0
total_params_reference = 1000.0
\end{lstlisting}

A model with \(N_{\mathrm{ref}}\) parameters is scored as \(2L\).
A model with half that many parameters is scored as \(1.5L\).
Choose the reference near the middle of the window you care about.
With the shipped reference of \(10^{5}\) and a few hundred parameters, the factor is about \(1.006\), so a multiplier of weight 1 barely moves the ranking.
The shipped eddy file avoids that ambiguity by setting the mode to \code{raw_loss}.
In \code{loss_multiplier} mode the weight is live.

To prefer the smallest model that reaches a loss target, first finish a classical baseline and read its \code{raw_objective}.
Then start a new study:

\begin{lstlisting}
optimize_total_parameters = True
parameter_objective_mode = "target_error_min_params"
target_raw_objective = 1.0e-3  # replace with the baseline raw loss
target_error_constraint_weight = 100.0
total_params_reference = 100.0
\end{lstlisting}

Trials above the target are dominated by the excess term.
Trials at or below the target are ranked by \(N/N_{\mathrm{ref}}\).
Use a new \code{study_name}.
The baseline trials were ranked by a different objective, and Optuna's best trial would mix the two scales if they shared a study.

\section{Parallel workers}

\file{run_cores.sh} runs the search as \(N\) processes and then trains the best trial once.

\begin{lstlisting}[style=sh]
./run_cores.sh 4
\end{lstlisting}

\(N\) is a positive integer.
The script creates \file{logs/}, deletes \file{logs/trial_table_rows} so the table header is printed once for this launch, and sets
\code{OMP_NUM_THREADS}, \code{MKL_NUM_THREADS}, \code{TF_NUM_INTRAOP_THREADS}, and \code{TF_NUM_INTEROP_THREADS} to 1.
Quantum trials run a compiled JAX circuit.
One thread per process keeps the workers from oversubscribing the machine.
The script does not pin JAX's own thread pool beyond those variables; a quantum worker can still use more than one core during compilation.

Each worker is \code{python3 run_pinn.py --search-only --worker-id i}, with standard output and standard error in \file{logs/worker_i.txt}.
Completed-trial rows are also written to the terminal.
The workers share the SQLite study.
The first worker creates the database under a file lock.
The storage layer waits up to 120 seconds on a locked database.
Together the workers stop near \code{n_trials} completed trials.
When the parameter window is set, discarded draws do not count, and the workers keep going until the cap of completed trials is reached.

If any worker exits with an error, the script does not start final training.
Read \file{logs/worker_*.txt}.
If every worker succeeds, the script prints \code{Search finished. Training the best completed trial.} and runs \code{python3 run_pinn.py --finalize}.

Do not start a second \code{--finalize}, or a plain \code{python run_pinn.py}, while the workers are still writing the study.
A plain invocation both searches and then trains.
Two final training jobs would write the same weight file.

You can split the same work by hand:

\begin{lstlisting}[style=sh]
python run_pinn.py --search-only --worker-id 0
python run_pinn.py --search-only --worker-id 1
python run_pinn.py --finalize
\end{lstlisting}

Run the two search commands in separate terminals, wait until both have returned, and then finalize.
\code{--search-only} without \code{--worker-id} is an error.

\section{Continuing, extending, and restarting}

The study is reopened by \code{study_name} inside \code{storage_url}.
Completed trials already in the file count toward \code{n_trials}.

\begin{itemize}
\item To add trials, raise \code{n_trials} and run again.
  Leave the family list and the integer boxes unchanged.
\item To change \code{model_family_options}, \code{embedding_options}, or the min, max, or step of a parameter that has already been sampled, create a new \code{study_name} and a new \code{storage_url}.
  Optuna stores the distribution with the study and rejects a conflicting suggestion.
\item To abandon a smoke-test study, point both names at a new database.
  You may delete the old \file{analysis/<name>/} folder when you no longer need its weights.
\item \code{--finalize} never adds trials.
  It trains the best completed trial under the current objective.
  If the study has no completed trial, it raises \code{RuntimeError} and names the study.
\end{itemize}

\section{Which trial is ``best''}

During final training the solver uses Optuna's best trial, the completed trial with the lowest stored objective.
That objective is \(L\) in \code{raw_loss} mode and the transformed value in the other modes.
The analysis script, described in the next chapter, selects the completed trial with the lowest \code{raw_objective} user attribute.
When the mode is \code{raw_loss}, the two selections are the same trial.
When the mode charges for parameter count, Optuna's best trial and the lowest-residual trial can differ.
The summary written by \code{--finalize} records Optuna's best trial.
The analytic report records the lowest residual.
Read the trial number in each file before you compare them.

\chapter{Outputs and the study report}
\label{ch:outputs}

\section{Files written by the trainer}

For a database \code{sqlite:///NAME.db} the trainer writes under \file{analysis/NAME/}.

\begin{description}
\item[\file{NAME.db}]
The Optuna study, present only when a search ran.
A fixed-architecture run does not create it.
SQLite may also leave \file{NAME.db-wal} and \file{NAME.db-shm} while a process has the file open.
A \file{NAME.db.lock} file serializes creation of a new database.

\item[\file{data/best_hyperparameters.txt}]
Family, width, depth, learning rate, PDE weights, and, for a quantum family, qubits, layers, and embedding.
Then the total, classical, and quantum trainable counts.
If a search ran, the Optuna trial number, its objective, and the stored loss pieces.
If L-BFGS-B ran, its success flag, message, iteration count, and evaluation count.
A skipped L-BFGS-B stage is recorded as skipped.

\item[\file{data/best_optuna_model.weights.h5}]
Keras weights of the final model.

\item[\file{data/best_optuna_model_config.json}]
The \code{HyperParams} record: width, depth, learning rate, both PDE weights, both log-weights, family, qubits, layers, and embedding.

\item[\file{data/<problem>_losses_data.npz}]
Adam total, mass, and momentum histories, and the same three series for L-BFGS-B when that stage ran.
The metadata repeats the family and the L-BFGS status.
\end{description}

Field archives depend on the benchmark.
One-dimensional problems write

\begin{itemize}
\item \file{<problem>_elevation_compare_data.npz} with arrays \code{TIME}, \code{X}, \code{A} (reference), and \code{M} (model);
\item \file{<problem>_velocity_compare_data.npz} with the same layout for \(u\);
\item \file{<problem>_error_data.npz} with \code{ETA_ERR} and \code{U_ERR};
\item \file{<problem>_profiles_boundary_data.npz} with elevation profiles and the velocity at the left endpoint.
\end{itemize}

The tide archives are converted to dimensional elevation, velocity, kilometers, and hours.
The dam break and the bump stay in the units of their dataclasses.
Steady square problems write \file{elevation}, \file{velocity_u}, and \file{velocity_v} compare files, plus \file{<problem>_error_data.npz} with \code{H_ERR}, \code{U_ERR}, and \code{V_ERR}.
Each archive carries a JSON string \code{meta_json} with the family, the sizes, and the constraint mode.

Load an archive from a short Python session:

\begin{lstlisting}
import json
import numpy as np

path = "analysis/optuna_manual_classical_dambreak/data/dam_break_elevation_compare_data.npz"
with np.load(path, allow_pickle=False) as z:
    model = z["M"]
    reference = z["A"]
    meta = json.loads(str(z["meta_json"]))
print(meta["model_family"], model.shape)
print(float(np.max(np.abs(model - reference))))
\end{lstlisting}

The trainer does not write PDF figures.
\file{swe_pinn/plots/jcp.py} contains helpers for loss curves and field panels.
The report script in the next section is the path that turns a finished study into figures.

\section{The study report}

From the repository root, with the database named in \file{user_input.py}:

\begin{lstlisting}[style=sh]
python analysis/analyze_study.py
\end{lstlisting}

The script reads that study, writes tables and figures under \file{analysis/<database name>/}, and compares the lowest-residual completed trial with the exact solution.
Override the locations when the active input file points somewhere else:

\begin{lstlisting}[style=sh]
python analysis/analyze_study.py \
    --db optuna_manual_eddies_three_family.db \
    --study manual_eddies_three_family \
    --problem steady_eddies
\end{lstlisting}

\code{--skip-train} writes the tradeoff report and the Optuna figures and skips the field comparison.
\code{--target} draws an optional horizontal line at that recorded objective in the text report.

\subsection{Text}

\file{data/param_tradeoff.txt} ranks completed trials by the recorded residual and lists family, size, and parameter counts.
That residual is the validation loss.
It is not the analytic field error.

\file{data/analytic_comparison.txt} is written when the comparison step runs.
It names the trial, the family, the architecture, the parameter counts, whether the weights were loaded or trained, the mesh, and for each field the discrete \(L_{2}\) norm of the error, the relative \(L_{2}\) error, and the \(L_{\infty}\) error.
\file{data/analytic_comparison.npz} stores the fields used for those numbers.

If saved weights match the selected architecture, they are loaded and not trained again.
If they are missing or were written for a different family or size, the script trains that architecture for \code{final_adam_epochs} on \code{final_adam_nx} by \code{final_adam_nt}, then saves the new weights.
A fixed-architecture example has no study, so this script has nothing to read until a search has completed at least one trial.

\subsection{Figures}

The script writes the following PDFs in \file{figures/} when the corresponding data exist.

\begin{center}
\begin{tabular}{@{}lp{8.6cm}@{}}
\toprule
File & Contents \\
\midrule
\file{optuna_objective_history.pdf} &
Objective against trial number \\
\file{optuna_error_vs_parameters.pdf} &
Residual against trainable-parameter count \\
\file{optuna_pareto_front.pdf} &
The nondominated residual-versus-size front \\
\file{optuna_model_family_comparison.pdf} &
Families side by side \\
\file{optuna_qubits_vs_error.pdf} &
Best residual at each qubit count \\
\file{optuna_quantum_layers_vs_error.pdf} &
Best residual at each layer count \\
\file{optuna_parameter_importance.pdf} &
fANOVA importance inside each family, one panel per family, each panel normalized on its own \\
\file{model_vs_analytic.pdf} &
Model, reference, and error for each field \\
\bottomrule
\end{tabular}
\end{center}

Family colors and markers are fixed in the script: classical is a blue circle, hybrid quantum is an orange square, and full quantum is a green triangle.
A reupload series is drawn separately when a study contains both embeddings.
A study that fixes \code{reupload}, as the shipped files do, is reported as the reupload series of that family.
Importance is skipped for a family with too few completed trials for fANOVA, and the script prints that it skipped the panel.

On \file{model_vs_analytic.pdf} each field is a row.
The columns are the model, the reference, and the signed error.
The horizontal axis is time for a two-field problem and \(y\) for a three-field problem.
The vertical axis is \(x\).
The error panel is annotated with the relative \(L_{2}\) error of that field.

\section{How to read a three-family figure}

On the error-versus-parameters figure, a family that sits lower at the same horizontal position reached a smaller residual at that size.
The parameter window in Section~\ref{sec:window} is what makes that horizontal comparison possible: draws outside the overlap were discarded, so the plotted points lie in a range every family was allowed to enter.
Full-quantum points still tend to lie toward the left of that window, because a pure circuit with a handful of qubits is smaller than a width-5, depth-7 perceptron.
Classical points can occupy the whole window.
Hybrid points occupy the part of the window that remains after the circuit is added to the trunk.

The Pareto figure is the set of completed trials that are not worse in both residual and size than some other completed trial.
A family can be absent from that front.
That is a result about the study you ran, at the epoch budget you set.
A 50-epoch smoke test does not support the same reading as a 4000-epoch study.

The analytic page is one model, the lowest residual in the study, retrained or reloaded at the final-stage budget.
Quote the trial number and the family from \file{analytic_comparison.txt} together with the relative \(L_{2}\) numbers.
The Optuna objective of that trial is a residual on the score lattice.
The relative \(L_{2}\) number is a field error on the evaluation mesh.
They answer different questions.

\chapter{Troubleshooting}
\label{ch:trouble}

\section{The family or the counts look wrong}

The banner is the authority.
Compare it with Table~\ref{tab:counts}.

\begin{itemize}
\item \code{model_family=classical} together with a positive \code{n_qubits_value} means the run is classical: either you set \code{model_family = "classical"}, or a hybrid specification had zero qubits and was rewritten.
  The rewrite is printed by the stored family on the trial, which becomes \code{classical}.
\item \code{width=0} and \code{depth=0} on a full-quantum banner are expected.
  Those fields do not enter Equation~\eqref{eq:full-count}.
\item A hybrid total far above the classical total at the same width and depth is the encoder, the decoder, and the angles.
  Equation~\eqref{eq:hybrid-count} is the check.
\item Quantum parameters count only the circuit angles, \(3 m n_{q} n_{L}\) for \(m\) fields.
  Encoder and decoder weights are classical, including on a full-quantum model.
\end{itemize}

\section{Optuna rejects the study}

A message about a dynamic value space, a categorical distribution, or an incompatible distribution means the new suggestion does not match trials already stored under \code{study_name}.
The usual edits that cause it are adding or removing a family, adding or removing an embedding, or changing the min, max, or step of width, depth, qubits, or layers.

Point \code{storage_url} at a new file and set a new \code{study_name}.
Do not reuse the old name inside the new file either, if you copy the database by hand.
Confirm the path the trainer will open: a relative URL \code{sqlite:///NAME.db} is \file{analysis/NAME/NAME.db}.

\section{The search finishes without training the model you expected}

If \code{n_trials} completed trials are already in the study, a new process adds nothing and trains the current best trial.
Raise the cap, or use a new study.

If \code{use_optuna} is false, the tune list is empty, or \code{n_trials} is 0, the log says the search was skipped.
The model that trains is the fixed \code{*_value} architecture.
That is the behavior of Chapter~\ref{ch:examples}.

\code{--search-only} never trains the final model.
Run \code{--finalize} after the workers return.
\code{--finalize} and \code{--search-only} together are an error, and \code{--search-only} without a worker id is an error.

\section{Every quantum trial is discarded}

The parameter window is in force, and the sampled circuits do not fit.
A full-quantum model with many qubits and layers exceeds \code{total_params_max}.
A classical model smaller than \code{total_params_min} is discarded as well.
Read the discard line: it prints the count, the family, and the bounds.
Either widen the window, using Table~\ref{tab:counts} to recompute the overlap, or shrink \code{n_qubits_max} and \code{q_layers_max}.
Setting both bounds to \code{None} disables the window.
Do that in a new study if the old study's interpretation depended on the overlap.

\section{The loss columns}

Hard mode prints IC and BC as zero.
A nonzero IC or BC value in hard mode means the constraint string was not \code{hard}.
The legal strings are exactly \code{hard} and \code{soft}, in either letter case.

Soft mode can show a small PDE loss and a large penalty.
The Optuna objective includes the penalty.
A flat field can make a steady residual small while missing the boundary data; the penalty is what prevents that trial from winning.
If the penalty dominates for the whole trial, the boundary weights on the physical dataclass are large relative to the residual, or the network has not yet fit the trace.
The dam-break weights of 50 on the initial step and 10 on the walls are the reference balance for that problem.

Equal log-weights produce PDE weights \((1,1)\).
If the mass column is systematically smaller than the momentum column, that is the residual, not a hidden extra factor, as long as both log-weights are zero.

\section{Adam prints nothing}

\code{show_epoch_output} or \code{show_trial_epoch_output} is false, or the stride is larger than the epoch budget.
The final loss is still stored in the loss archive.
Trial rows are a separate print and still appear when a trial completes.

\section{Workers fail together}

\begin{itemize}
\item \code{One or more workers failed} from \file{run_cores.sh} means at least one \file{logs/worker_*.txt} ends in a traceback.
  Final training was not started.
\item \code{database is locked} that persists past the 120-second wait usually means a crashed process left a transaction, or two launches are sharing the file.
  Stop the extra launch.
  If no process has the database open, remove a stale \file{-wal} only when you have confirmed that no writer is alive.
  The trainer relocates a legacy database from the repository root or from \file{analysis/<problem name>/} into \file{analysis/<database name>/} on open.
  Two copies of the same file in those old locations are not merged; an existing destination is left in place.
\item A worker that cannot import PennyLane fails on its first quantum trial.
  A classical trial in the same study still completes.
  Install PennyLane and rerun; completed trials are kept.
\item An out-of-memory failure on a quantum trial is a large collocation grid times a large circuit.
  The circuit is evaluated as a statevector, batched over the collocation points.
  Reduce \code{trial_nx} and \code{trial_nt}, or reduce \code{n_qubits_max}, before raising the worker count.
\end{itemize}

\section{The analytic comparison retrains}

The script prints that it is training the best trial when the weight file is missing or does not match the architecture.
Matching requires the same family, sizes, learning rate, weights, and embedding.
A smoke test that overwrote \file{best_optuna_model.weights.h5} in a production folder is the usual cause.
Restore the weights from your own copy, or let the script train, which replaces the file with the architecture it selected.
Keep smoke tests on a separate \code{storage_url}, as in Chapter~\ref{ch:examples}.

The comparison selects the lowest \code{raw_objective}.
Final training selects Optuna's best objective.
Section on which trial is best, in Chapter~\ref{ch:search}, is the distinction.
Quote the trial number from the file you are reading.

\section{Full quantum is slow, and warmup did nothing}

Compilation happens on the first forward pass of each distinct circuit size in the process.
A search that draws many \((n_{q}, n_{L})\) pairs compiles many circuits.
That cost is paid again in every worker.

\code{hybrid_classical_warmup_fraction} does not apply to \code{full_quantum}.
The circuit is the model.
To keep a hybrid circuit off for the first quarter of a trial, set the fraction to \code{0.25} and confirm the warmup line in the trial log.
The score is always computed with the circuit on.

\section{Checklist before a long study}

\begin{enumerate}
\item The banner of a one-trial or fixed run shows the family, the embedding, and the parameter counts you computed.
\item \code{problem_name}, \code{constraints}, and the physical dataclass are the benchmark you mean to rank.
\item \code{storage_url} and \code{study_name} are new, or you intend to append to that study.
\item \code{model_family_options} lists only the families that should appear.
  A two-family preset is a preset that already removed \code{full_quantum}.
\item The parameter window matches the number of fields.
  The numbers 42 and 603 belong to the three-field eddy box.
\item \code{optuna_n_startup_trials} is above \code{n_trials} when you want equal family draws, and below it when you want TPE to take over.
\item \code{trial_epochs} is the budget you are willing to rank, and \code{final_adam_epochs} is the budget for the model you will plot.
\item \code{use_lbfgs} is the choice you want recorded in the summary.
\item PennyLane imports in a one-line hybrid or full-quantum build before you occupy a queue with hundreds of trials.
\end{enumerate}


\end{document}
